%%=========================================================================%%
%%                               Settings                                  %%
%%=========================================================================%%

% Uncomment for PhD mode (default is licentiate mode)
% \def\phdThesis{1}

\newcommand{\currentyear}{2026}
\newcommand{\authorname}{AUTHOR NAME}
\newcommand{\mytitle}{THESIS TITLE}
\newcommand{\mysubtitle}{Optional Subtitle}

\newcommand{\division}{Division or Research Area}
% \newcommand{\researchgroup}{Optional Research Group}

% PhD-specific identifiers
\newcommand{\phdISBN}{ISBN-PHD-PLACEHOLDER}
\newcommand{\phdSeriesNumber}{SERIES-NUMBER}
\newcommand{\techReportNumber}{TECH-REPORT-NUMBER}

\newcommand{\licISSN}{LIC-ISSN}
\newcommand{\phdISSN}{PHD-ISSN}

\newcommand{\chalmers}{University Name}
\newcommand{\gotuni}{Partner Institution}
\newcommand{\mydepartment}{Department Name}
\newcommand{\chalIgu}{\chalmers}
\newcommand{\chalandgu}{\chalmers~and~\gotuni}

\ifx\phdThesis\undefined
\newcommand{\degreetitle}{Licentiate of Discipline}
\else
\newcommand{\degreetitle}{Doctor of Philosophy}
\fi

\ifx\phdThesis\undefined
\newcommand{\identifierNoText}{
ISSN \licISSN\\
}
\else
\newcommand{\identifierNoText}{ISBN \phdISBN\\
Series number \phdSeriesNumber\\
ISSN \phdISSN\\
Technical Report No. \techReportNumber \\
\vspace{1.5cm}
}
\fi

\setlength{\headheight}{14mm}
\usepackage{geometry}

\ifx\inlagaPage\undefined
\geometry{paperwidth=169mm, paperheight=239mm, inner=28mm,outer=22mm,top=22mm, bottom=18mm}
\else
\geometry{paperwidth=148mm, paperheight=210mm, inner=25mm,outer=25mm,top=22mm, bottom=18mm}
\fi

\usepackage{fancyhdr}
\pagestyle{fancy}
\fancyhead[LE,RO]{\thepage}
\fancyhead[RE]{\scriptsize \slshape \leftmark}
\fancyhead[LO]{\scriptsize \slshape \rightmark}
\fancyfoot[C]{}

\setcounter{secnumdepth}{3} \setcounter{tocdepth}{3}

\usepackage[labelformat=simple]{subcaption}
\renewcommand\thesubfigure{(\alph{subfigure})}
\usepackage{graphicx}
\usepackage{url}
\usepackage[british]{babel}
\usepackage[backend=bibtex,style=ieee,maxcitenames=3,backref=true,language=british]{biblatex}
\addbibresource{../refs/references.bib}

\usepackage{amsmath,amssymb,amscd,latexsym,dsfont}
\usepackage{textcomp}
\usepackage{psfrag}
\usepackage{rotating}
\usepackage{enumitem}
\usepackage{pdflscape}
\usepackage{tikz}
\usetikzlibrary{positioning,arrows.meta,fit,calc,shapes.geometric}
\usepackage{adjustbox}
\usepackage{booktabs,tabularx,array,longtable}
\usepackage{ragged2e}
\usepackage[acronym, nopostdot, nonumberlist]{glossaries-extra}
% Template glossary placeholders
\newacronym{api}{API}{Application Programming Interface}
\newglossaryentry{thesis}{
  name={thesis},
  description={A structured scholarly document presenting research work}
}

\makeglossaries
\setglossarystyle{list}
\setglossarysection{section}

\usepackage[final]{pdfpages}
\usepackage{xspace}
\usepackage{relsize}
\usepackage{IEEEtrantools}
\usepackage{microtype}
\usepackage[textsize=tiny,textwidth=20mm,backgroundcolor=orange!70]{todonotes}
\setlength{\marginparwidth}{20mm}

\newcommand{\figref}[1]{Fig.~\ref{#1}}
\newcommand{\tabref}[1]{Table~\ref{#1}}
\newcommand{\chapref}[1]{Chap.~\ref{#1}}
\newcommand{\secref}[1]{Section~\ref{#1}}
\newcommand{\eqnref}[1]{Equation~\ref{#1}}
\newcommand{\paperref}[1]{Paper~\romannum{#1}}

\newcommand{\fe}[1]{\textsf{\small #1}}
\newcommand{\p}[1]{\textsf{\small #1}}

\newcommand{\interviewquote}[2]{\begin{quote}
\footnotesize{\emph{``#1'' }} --- \footnotesize{#2}
\end{quote} }

\makeatletter
\newcommand*{\IfItalicsTF}{%
  \ifx\f@shape\my@test@it
    \expandafter\@firstoftwo
  \else
    \expandafter\@secondoftwo
  \fi
}
\newcommand*{\my@test@it}{it}
\makeatother

\newcommand{\romannum}[1]{{\uppercase\expandafter{\romannumeral #1\relax}}}
\newcommand{\scalefactorOO}{0.45}
